\chapter{Complexity — TMSAT}
%%% generated by the blueprint pipeline from TCSlib.Complexity.ClassNP.TMSAT
%%%%%%%%%%%%%%%%%%%%%%%%%%%%%%%%%%%%%%%%%%%%%%%%%%%%%%%%%%%%%%%%%%%%%%%%%%%%%%%
\section{Overview}

This module defines the language $\mathrm{TMSAT}$ of [AB09, Theorem~2.9]: over a
machine-representation scheme $c$, the right-nested pair encodings
$\langle\alpha,\langle x,\langle 1^n,1^t\rangle\rangle\rangle$ such that the machine coded by
$\alpha$ outputs $1$ on $\langle x,u\rangle$ within $t$ steps for some certificate $u$ of length
$n$. It proves that $\mathrm{TMSAT}$ is $\mathrm{NP}$-hard for every scheme and lies in
$\mathrm{NP}$ (hence is $\mathrm{NP}$-complete) for effective schemes whose canonizer runs in
polynomial time; along the way it shows that $n \mapsto C(n+1)^{c+1}$ is time constructible for
$C > 0$ via an explicit unary box-enumeration machine, gives a quantitative timed universal
simulation bound, and builds the polynomial-time parsing, prefix-extraction and request-assembly
layer of the verifier.

%%%%%%%%%%%%%%%%%%%%%%%%%%%%%%%%%%%%%%%%%%%%%%%%%%%%%%%%%%%%%%%%%%%%%%%%%%%%%%%
\section{Declarations}

\begin{definition}[The language TMSAT]
\lean{Complexity.TMSAT}
\label{Complexity.TMSAT}
\leanok
\uses{Turing.CodeTM.toFinTM, Turing.FinTM.ComputesInTime, Turing.MachineCode, Turing.pairEncode}
Let $c$ be a machine-representation scheme, which assigns to every bit string $\alpha$ a one-work-tape binary Turing machine $M_\alpha$ (decoding is total). The language $\mathrm{TMSAT}_c$ consists of all bit strings $\langle\alpha,\langle x,\langle 1^n,1^t\rangle\rangle\rangle$, where $\langle a,b\rangle$ is the self-delimiting pair encoding (every bit of $a$ doubled, then $01$, then $b$) and $1^k$ is the word of $k$ ones, for which there exists a certificate $u$ of length exactly $n$ such that $M_\alpha$, started on input $\langle x,u\rangle$, halts within $t$ steps with output exactly the single bit $1$.
\end{definition}

\begin{definition}[Control states of the unary box enumerator]
\lean{Complexity.PolyControl}
\label{Complexity.PolyControl}
\leanok
For natural numbers $c$ and $C$, this is the finite set of control states of the unary box-enumeration machine: a copy state, a setup state, a loop state $\mathrm{loop}_i$ and a rewind state $\mathrm{rewind}_i$ for each level $i \in \{0,\dots,c\}$, an advance state $\mathrm{advance}_i$ for each $i \in \{0,\dots,c+1\}$, and an emission state $\mathrm{emit}_j$ for each $j \in \{0,\dots,C\}$, recording how many output bits remain to be emitted.
\end{definition}

\begin{definition}[Unary tape contents]
\lean{Complexity.polyTape}
\label{Complexity.polyTape}
\leanok
For $q \in \mathbb{N}$, the work-tape contents $\tau_q \colon \mathbb{Z} \to \{0,1,\sqcup\}$ holding the unary word $1^q$ are given by $\tau_q(z) = 1$ if $0 \le z < q$, and $\tau_q(z) = \sqcup$ (blank) otherwise.
\end{definition}

\begin{definition}[Single-head move action]
\lean{Complexity.polyMove}
\label{Complexity.polyMove}
\leanok
\uses{Complexity.PolyControl, Turing.Action}
For $c, C \in \mathbb{N}$, a work-tape index $i \in \{0,\dots,c\}$, a direction $d \in \{-1,0,1\}$ and a control state $s$ of the unary box enumerator, this is the one-step action of a $(c+1)$-work-tape binary machine that keeps the input head in place, writes nothing on any work tape, moves work head $i$ by $d$ while leaving every other work head in place, emits no output, and passes to state $s$.
\end{definition}

\begin{definition}[Unary box-enumeration machine]
\lean{Complexity.polyUnaryTM}
\label{Complexity.polyUnaryTM}
\leanok
\uses{Complexity.PolyControl, Complexity.polyMove, Turing.FinTM}
For $c, C \in \mathbb{N}$, $\mathcal{P}_{c,C}$ is the finite binary-alphabet machine with $c+1$ work tapes, control states copy, setup, $\mathrm{loop}_i$, $\mathrm{rewind}_i$ ($0 \le i \le c$), $\mathrm{advance}_i$ ($0 \le i \le c+1$) and $\mathrm{emit}_j$ ($0 \le j \le C$), and initial state copy. In copy it scans the input, writing $1$ on every work tape and moving all heads right; at the end of the input it writes one more $1$ everywhere, moves all heads left and enters setup, which moves all work heads left until head $0$ reads a blank, then one cell right into $\mathrm{loop}_c$. In $\mathrm{loop}_i$, if head $i$ reads a blank it moves head $i$ left into $\mathrm{rewind}_i$, and otherwise it passes to $\mathrm{loop}_{i-1}$, or to $\mathrm{emit}_C$ when $i = 0$; $\mathrm{emit}_j$ for $j > 0$ outputs a $1$ and passes to $\mathrm{emit}_{j-1}$, and $\mathrm{emit}_0$ passes to $\mathrm{advance}_0$; $\mathrm{rewind}_i$ moves head $i$ left until it reads a blank, then one cell right into $\mathrm{advance}_{i+1}$; $\mathrm{advance}_i$ for $i \le c$ moves head $i$ one cell right into $\mathrm{loop}_i$, and $\mathrm{advance}_{c+1}$ halts.
\end{definition}

\begin{definition}[Loop configurations of the box enumerator]
\lean{Complexity.polyCfg}
\label{Complexity.polyCfg}
\leanok
\uses{Complexity.PolyControl, Complexity.polyTape, Turing.Cfg}
Let $c, C, q \in \mathbb{N}$ and let $x$ be a bit string. For a control state $s$ of the unary box-enumeration machine $\mathcal{P}_{c,C}$, head positions $h \colon \{0,\dots,c\} \to \mathbb{Z}$ and an output word $o$, the loop configuration $\Gamma_q(s,h,o)$ is the configuration of a $(c+1)$-work-tape machine on input $x$ that is in state $s$, has its input head on the blank cell immediately right of the input, has every work tape holding $1^q$ in cells $0,\dots,q-1$ and blanks elsewhere, has work head $k$ at position $h(k)$, and has output $o$.
\end{definition}

\begin{lemma}[Effect of a Single-Head Move]
\lean{Complexity.polyMove_apply}
\label{Complexity.polyMove_apply}
\leanok
\uses{Complexity.PolyControl, Complexity.polyCfg, Complexity.polyMove, Turing.Action.apply, Turing.moveInputPos_zero}
\difficulty{3}
Let $c,C,q\in\mathbb{N}$, let $x$ be a finite bit string, let $s,s'$ be control states,
let $h:\{0,\ldots,c\}\to\mathbb{Z}$ be a list of work-head positions, and let $o$ be an
output bit string. For an index $i\in\{0,\ldots,c\}$ and a direction $d\in\{-1,0,1\}$,
applying the action that moves work head $i$ by $d$ and enters state $s'$ to the
configuration with state $s$, head positions $h$, and output $o$ yields the
configuration with state $s'$, output $o$, and head positions obtained by replacing
$h(i)$ by $h(i)+d$.
\end{lemma}

\begin{lemma}[Emission chain appends its remaining bits]
\lean{Complexity.poly_emit}
\label{Complexity.poly_emit}
\leanok
\uses{Complexity.polyCfg, Complexity.polyUnaryTM, Turing.Action.apply, Turing.MultiTapeTM.runFrom, Turing.MultiTapeTM.runFrom_succ_eq_step, Turing.MultiTapeTM.runFrom_zero, Turing.MultiTapeTM.step}
\difficulty{4}
Let $c,C,q\in\mathbb{N}$, let $x$ be a bit string, let $h:\{0,\ldots,c\}\to\mathbb{Z}$
be a head-position function, and let $j\le C$ and $o$ be an output word. After exactly
$j+1$ steps of the unary box-enumeration machine from the configuration on input $x$ in
state $\mathrm{emit}_j$, with tape contents $1^q$ and head positions $h$, the machine is
in state $\mathrm{advance}_0$, with the same tape contents and head positions, and
output $o$ followed by $j$ copies of $1$.
\end{lemma}

\begin{lemma}[Rewinding a unary loop head]
\lean{Complexity.poly_rewind}
\label{Complexity.poly_rewind}
\leanok
\uses{Complexity.PolyControl, Complexity.polyCfg, Complexity.polyMove_apply, Complexity.polyTape, Complexity.polyUnaryTM, Turing.Action, Turing.Action.apply, Turing.Cfg.workTapeSymbols, Turing.MultiTapeTM.runFrom, Turing.MultiTapeTM.runFrom_succ_eq_step, Turing.MultiTapeTM.runFrom_zero, Turing.MultiTapeTM.step}
\difficulty{4}
Let $c,C,q\in\mathbb{N}$, let $x$ be a bit string, let $h:\{0,\ldots,c\}\to\mathbb{Z}$
be a head-position function, let $o$ be an output word, and let $i\in\{0,\ldots,c\}$.
For every $j\le q$, after exactly $j+1$ steps of the unary box-enumeration machine from
the configuration on input $x$ in state $\mathrm{rewind}_i$, with tape contents $1^q$
and head positions equal to $h$ except that head $i$ is at $j-1$, it reaches state
$\mathrm{advance}_{i+1}$ with head positions equal to $h$ except that head $i$ is at
$0$, and with output $o$.
\end{lemma}

\begin{lemma}[Advancing a loop head]
\lean{Complexity.poly_advance}
\label{Complexity.poly_advance}
\leanok
\uses{Complexity.polyCfg, Complexity.polyMove_apply, Complexity.polyUnaryTM, Turing.MultiTapeTM.step}
\difficulty{2}
Let $c,C,q\in\mathbb{N}$, let $x$ be a bit string, let $h:\{0,\ldots,c\}\to\mathbb{Z}$
be a head-position function, let $o$ be an output word, and let $i\in\{0,\ldots,c\}$.
One step of the unary box-enumeration machine from the configuration on input $x$ in
state $\mathrm{advance}_i$, with tape contents $1^q$ and head positions $h$, reaches the
configuration in state $\mathrm{loop}_i$ with the same tape contents and output, and
with head positions equal to $h$ except that head $i$ is at $h(i)+1$.
\end{lemma}

\begin{definition}[Running time of nested unary loops]
\lean{Complexity.polyCost}
\label{Complexity.polyCost}
\leanok
For $q, C \in \mathbb{N}$, the function $\kappa_{q,C} \colon \mathbb{N} \to \mathbb{N}$ is defined by $\kappa_{q,C}(0) = C + 1$ and $\kappa_{q,C}(r+1) = q\,(\kappa_{q,C}(r) + 2) + q + 2$; it is the exact number of steps of a full nest of $r$ unary loops of side $q$ that emit $C$ bits per innermost iteration.
\end{definition}

\begin{lemma}[Correctness and cost of a unary loop]
\lean{Complexity.poly_loop}
\label{Complexity.poly_loop}
\leanok
\uses{Complexity.polyCfg, Complexity.polyCost, Complexity.polyMove_apply, Complexity.polyTape, Complexity.polyUnaryTM, Complexity.poly_advance, Complexity.poly_emit, Complexity.poly_rewind, Turing.Action.apply, Turing.Cfg.workTapeSymbols, Turing.MultiTapeTM.runFrom, Turing.MultiTapeTM.runFrom_add, Turing.MultiTapeTM.runFrom_succ_eq_step, Turing.MultiTapeTM.runFrom_zero, Turing.MultiTapeTM.step}
\difficulty{6}
Let $c,C,q\in\mathbb{N}$ with $q>0$, let $x$ be a bit string, let $i<c+1$, and let
$h:\{0,\ldots,c\}\to\mathbb{Z}$ satisfy $h(k)=0$ whenever $k\le i$. Let $o$ be an output
word, and let $r,j\in\mathbb{N}$ satisfy $j+r=q$. After exactly
$r(\kappa_{q,C}(i)+2)+q+2$ steps of the unary box-enumeration machine from the
configuration on input $x$ in state $\mathrm{loop}_i$, with tape contents $1^q$, output
$o$, and head positions equal to $h$ except that head $i$ is at $j$, it reaches state
$\mathrm{advance}_{i+1}$ with head positions $h$ and output $o$ followed by $rCq^i$
copies of $1$. Here $\kappa_{q,C}(0)=C+1$ and
$\kappa_{q,C}(m+1)=q(\kappa_{q,C}(m)+2)+q+2$.
\end{lemma}

\begin{lemma}[Extending a Unary Tape]
\lean{Complexity.polyTape_write}
\label{Complexity.polyTape_write}
\leanok
\uses{Complexity.polyTape}
\difficulty{3}
For every $q\in\mathbb{N}$, changing the tape contents $\tau_q$ by writing $1$ at cell
$q$ gives $\tau_{q+1}$.
\end{lemma}

\begin{lemma}[Bound on Nested-Loop Cost]
\lean{Complexity.polyCost_le}
\label{Complexity.polyCost_le}
\leanok
\uses{Complexity.polyCost}
\difficulty{4}
Let $q,C\in\mathbb{N}$ with $q>0$, and let $\kappa_{q,C}$ be defined by
$\kappa_{q,C}(0)=C+1$ and $\kappa_{q,C}(r+1)=q(\kappa_{q,C}(r)+2)+q+2$. Then, for every
$r\in\mathbb{N}$, $\kappa_{q,C}(r)\le (C+1+5r)q^r$.
\end{lemma}

\begin{definition}[Configurations of the copying phase]
\lean{Complexity.polyCopyCfg}
\label{Complexity.polyCopyCfg}
\leanok
\uses{Complexity.PolyControl, Complexity.polyTape, Turing.Cfg}
Let $c, C \in \mathbb{N}$, $x$ a bit string and $i \le |x|$. The copying configuration at stage $i$ is the configuration of the unary box-enumeration machine $\mathcal{P}_{c,C}$ (with $c+1$ work tapes) on input $x$ in the copy state, with the input head on input position $i$ (counting from $0$; the blank cell right of the input when $i = |x|$), every work tape holding $1^i$ in cells $0,\dots,i-1$ and blanks elsewhere, every work head at cell $i$, and empty output.
\end{definition}

\begin{lemma}[The input-copying phase]
\lean{Complexity.poly_copy}
\label{Complexity.poly_copy}
\leanok
\uses{Complexity.polyCopyCfg, Complexity.polyTape, Complexity.polyTape_write, Complexity.polyUnaryTM, Turing.Action.apply, Turing.Cfg.init, Turing.Cfg.inputSymbol, Turing.MultiTapeTM.initCfg, Turing.MultiTapeTM.runFrom, Turing.MultiTapeTM.runFrom_succ_eq_step', Turing.MultiTapeTM.step, Turing.inputSymbolInner, Turing.moveInputPos, Turing.moveInputPos_pos_of_ne_right}
\difficulty{4}
Let $c,C\in\mathbb{N}$ and let $x$ be a bit string. For every $i\le |x|$, after $i$
steps of the unary box-enumeration machine from its initial configuration on $x$, the
machine is in the copy state with its input head at position $i$ (counting from $0$),
each work tape containing $1^i$ in cells $0,\ldots,i-1$ and blanks elsewhere, every work
head at position $i$, and empty output.
\end{lemma}

\begin{lemma}[The setup rewind]
\lean{Complexity.poly_setup}
\label{Complexity.poly_setup}
\leanok
\uses{Complexity.polyCfg, Complexity.polyTape, Complexity.polyUnaryTM, Turing.Action.apply, Turing.Cfg.workTapeSymbols, Turing.MultiTapeTM.runFrom, Turing.MultiTapeTM.runFrom_succ_eq_step, Turing.MultiTapeTM.runFrom_zero, Turing.MultiTapeTM.step}
\difficulty{4}
Let $c,C,q\in\mathbb{N}$ and let $x$ be a bit string. For every $j\le q$, after exactly
$j+1$ steps of the unary box-enumeration machine from the configuration on input $x$ in
the setup state, with every work tape containing $1^q$, every work head at position
$j-1$, and empty output, it reaches the loop state at level $c$, with every work head at
position $0$, the same tape contents, and empty output.
\end{lemma}

\begin{lemma}[Startup of the box enumerator]
\lean{Complexity.poly_start}
\label{Complexity.poly_start}
\leanok
\uses{Complexity.polyCfg, Complexity.polyCopyCfg, Complexity.polyTape_write, Complexity.polyUnaryTM, Complexity.poly_copy, Complexity.poly_setup, Turing.Action.apply, Turing.Cfg.inputSymbol, Turing.MultiTapeTM.initCfg, Turing.MultiTapeTM.runFrom, Turing.MultiTapeTM.runFrom_add, Turing.MultiTapeTM.runFrom_succ_eq_step', Turing.MultiTapeTM.step, Turing.moveInputPos_zero}
\difficulty{4}
Let $c,C\in\mathbb{N}$ and let $x$ be a bit string. After exactly $2(|x|+1)$ steps from
its initial configuration on $x$, the unary box-enumeration machine is in the loop state
at level $c$, with its input head on the blank cell immediately to the right of the
input, every work tape containing $1^{|x|+1}$ in cells $0,\ldots,|x|$ and blanks
elsewhere, every work head at position $0$, and empty output.
\end{lemma}

\begin{lemma}[Unary polynomial computation]
\lean{Complexity.poly_unary_computes}
\label{Complexity.poly_unary_computes}
\leanok
\uses{Complexity.polyCfg, Complexity.polyCost, Complexity.polyCost_le, Complexity.polyUnaryTM, Complexity.poly_loop, Complexity.poly_start, Turing.Action.apply, Turing.FinTM.ComputesFunInTime, Turing.FinTM.ComputesInTime, Turing.FinTM.ComputesInTime.mono, Turing.FinTM.computesInTime_iff, Turing.MultiTapeTM.runFrom, Turing.MultiTapeTM.runFrom_add, Turing.MultiTapeTM.runFrom_succ_eq_step', Turing.MultiTapeTM.step}
\difficulty{4}
For all $c,C\in\mathbb{N}$, the unary box-enumeration machine with parameters $c,C$
computes the function sending each binary string $x$ to the word consisting of
$C(|x|+1)^{c+1}$ ones, in time $(C+5(c+1)+4)(|x|+1)^{c+1}$.
\end{lemma}

\begin{theorem}[Polynomial time constructibility]
\lean{Complexity.timeConstructible_poly}
\label{Complexity.timeConstructible_poly}
\leanok
\uses{Complexity.TimeConstructible, Complexity.poly_unary_computes, Complexity.timeConstructible_id, Turing.FinTM.ComputesFunInTime, Turing.FinTM.ComputesInTime.mono, Turing.FinTM.computesFunInTime_comp}
\difficulty{5}
For natural numbers $C$ and $c$ with $C>0$, the function $T(n)=C(n+1)^{c+1}$ is time
constructible: for every $n$, $n\le T(n)$, and some positive integer constant $a$ and
finite binary Turing machine compute, on every binary input $x$, the little-endian
binary representation of $T(|x|)$ within $a(T(|x|)+1)$ steps.
\end{theorem}

\begin{lemma}[Serialization Length Bound]
\lean{Complexity.tmsat_serialization_length}
\label{Complexity.tmsat_serialization_length}
\leanok
\uses{Turing.CodeTM.serialize, Turing.EffectiveMachineCode, Turing.FinTM.computesInTime_iff, Turing.MultiTapeTM.output_length_le}
\difficulty{2}
Let $c$ be an effective machine code with canonizer and let $\alpha$ be a binary word.
The length of the serialization of the machine decoded from $\alpha$ is at most the
canonizer's time bound at input length $|\alpha|$.
\end{lemma}

\begin{lemma}[Length under list flattening]
\lean{Complexity.tmsat_flatMap_length}
\label{Complexity.tmsat_flatMap_length}
\leanok
\difficulty{3}
Let $A$ and $B$ be types, let $f$ assign to each $a\in A$ a list over $B$ of length at
least $1$, and let $l$ be any finite list over $A$. Then the length of $l$ is at most
the length of the concatenation of the lists $f(a)$ as $a$ ranges through $l$.
\end{lemma}

\begin{lemma}[Nonempty action serialization]
\lean{Complexity.tmsat_action_nonempty}
\label{Complexity.tmsat_action_nonempty}
\leanok
\uses{Turing.Action, Turing.actionBits, Turing.optBoolBits, Turing.optOptBoolBits, Turing.optStateBits, Turing.signBits}
\difficulty{3}
For every natural number $n$ and every action of a one-work-tape binary Turing machine
with state set $\\{0,\\dots,n\\}$, the serialization of the action has length at least
$1$.
\end{lemma}

\begin{lemma}[Bounds from machine serialization]
\lean{Complexity.tmsat_serialization_parameters}
\label{Complexity.tmsat_serialization_parameters}
\leanok
\uses{Complexity.tmsat_action_nonempty, Complexity.tmsat_flatMap_length, Turing.CodeTM, Turing.CodeTM.serialize, Turing.actionBits, Turing.pairEncode, Turing.universal_pair_length}
\difficulty{5}
Let $M$ be a one-work-tape binary Turing machine in code normal form, with state set
$\{0,\ldots,m\}$, initial state $q_0$, and canonical serialization of length $L$. Then
the binary representation of $m$ has length at most $L$, the index of $q_0$ is at most
$L$, and $m+1\le L$.
\end{lemma}

\begin{lemma}[Concrete Simulation Coefficient Bound]
\lean{Complexity.tmsat_concrete_coefficient}
\label{Complexity.tmsat_concrete_coefficient}
\leanok
\uses{Complexity.tmsat_serialization_length, Complexity.tmsat_serialization_parameters, Turing.CodeTM.serialize, Turing.EffectiveMachineCode, Turing.universalBlockBound}
\difficulty{3}
Let $c$ be an effective machine code with canonizer and let $\alpha$ be a binary word.
The following inequality holds, where $L$ is the length of the serialization of the
machine decoded from $\alpha$, $m$ is its number of states minus one, $q_0$ is its
initial state index, and $T_c(n)$ is the canonizer's time bound at input length $n$:
\[
3|\alpha|+T_c(|\alpha|)+L+2|\mathrm{bin}(m)|+2q_0+16+\bigl(3L+5(m+1)+20\bigr)+14\le
3|\alpha|+14T_c(|\alpha|)+50.
\]
Here $\mathrm{bin}(m)$ is the little-endian binary representation of $m$.
\end{lemma}

\begin{theorem}[Quantitative Timed Universal Simulation]
\lean{Complexity.timed_universal_quantitative}
\label{Complexity.timed_universal_quantitative}
\leanok
\uses{Complexity.tmsat_concrete_coefficient, Turing.CodeTM.serialize, Turing.CodeTM.toFinTM, Turing.EffectiveMachineCode, Turing.FinTM, Turing.FinTM.ComputesInTime, Turing.FinTM.ComputesInTime.mono, Turing.pairEncode, Turing.timed_universal_concrete, Turing.universalBlockBound}
\difficulty{3}
Let $c$ be an effective machine code with canonizer, and let $\alpha,x$ be binary words
and $t\in\mathbb{N}$. There exists a finite-state binary Turing machine $U$ such that,
for every output word $y$, if the machine decoded from $\alpha$ halts on $x$ at exactly
time $t$ with output $y$, then $U$ halts on input
$\langle\langle\mathrm{bin}(t),\alpha\rangle,x\rangle$ at exactly time
$(3|\alpha|+14T_c(|\alpha|)+50)(t+1)^2$ with output $1y$; and if it does not halt on $x$
with any output at exactly time $t$, then $U$ halts on that same input at exactly that
time with output $0$. Here $T_c(n)$ is the canonizer's time bound at input length $n$,
$\mathrm{bin}(t)$ is the little-endian binary representation, and
$\langle\cdot,\cdot\rangle$ is the specified pairing of binary words.
\end{theorem}

\begin{definition}[Answer word of a bounded run]
\lean{Complexity.tmsatAnswer}
\label{Complexity.tmsatAnswer}
\leanok
\uses{Turing.MachineCode, Turing.MultiTapeTM.initCfg, Turing.MultiTapeTM.runFrom}
Let $c$ be a machine-representation scheme, which assigns to every bit string $\alpha$ a one-work-tape binary Turing machine $M_\alpha$ (decoding is total). For bit strings $\alpha, x$ and $t \in \mathbb{N}$, the answer word $A_c(\alpha, x, t)$ is $1w$ if $M_\alpha$, run for $t$ steps on input $x$, has halted with output $w$, and is the single bit $0$ if it has not halted after $t$ steps.
\end{definition}

\begin{lemma}[Timed Simulator Returns the Answer]
\lean{Complexity.tmsat_simulator_total}
\label{Complexity.tmsat_simulator_total}
\leanok
\uses{Complexity.tmsatAnswer, Turing.CodeTM.toFinTM, Turing.EffectiveMachineCode, Turing.FinTM, Turing.FinTM.ComputesInTime, Turing.FinTM.computesInTime_iff, Turing.MultiTapeTM.initCfg, Turing.MultiTapeTM.runFrom, Turing.pairEncode}
\difficulty{4}
Let $c$ be an effective machine-code scheme and $U$ a finite binary Turing machine.
Suppose that for every pair of binary strings $\alpha,x$ and every $t\in\mathbb{N}$, if
the machine decoded from $\alpha$ computes output $w$ on input $x$ in exactly $t$ steps,
then $U$ computes $1w$ on input $\langle\langle\mathrm{bin}(t),\alpha\rangle,x\rangle$
in exactly $S$ steps; and if it computes no output on $x$ in exactly $t$ steps, then $U$
computes the one-bit string $0$ on that encoded input in exactly $S$ steps, where
$S=(3|\alpha|+14T_c(|\alpha|)+50)(t+1)^2$. Then $U$ computes on that encoded input, in
exactly $S$ steps, the word that is $1w$ when the decoded machine has halted by time $t$
with output $w$, and is the one-bit string $0$ otherwise.
\end{lemma}

\begin{lemma}[Acceptance characterized by the bounded-run answer]
\lean{Complexity.tmsatAnswer_accept}
\label{Complexity.tmsatAnswer_accept}
\leanok
\uses{Complexity.tmsatAnswer, Turing.CodeTM.toFinTM, Turing.FinTM.ComputesInTime, Turing.FinTM.computesInTime_iff, Turing.MachineCode}
\difficulty{3}
Let $c$ be a machine-representation scheme, and let $\alpha,x$ be binary strings and
$t\in\mathbb{N}$. Run the machine decoded from $\alpha$ on $x$ for $t$ steps, and form
the answer word by prefixing its output with $1$ if it has halted by then, or by taking
the one-bit word $0$ if it has not halted. This answer word is $11$ if and only if the
decoded machine computes the one-bit output $1$ from $x$ within $t$ steps.
\end{lemma}

\begin{lemma}[Polynomial Simulation Budget]
\lean{Complexity.tmsat_simulation_budget}
\label{Complexity.tmsat_simulation_budget}
\leanok
\uses{Complexity.PolyBound, Turing.EffectiveMachineCode}
\difficulty{4}
Let $c$ be an effective machine-code scheme whose canonizer time function is
polynomially bounded: there exist $C,k\in\mathbb{N}$ such that $c$'s canonizer time at
input length $n$ is at most $C(n+1)^k$ for every $n\in\mathbb{N}$. Then there exist
$A,d\in\mathbb{N}$ such that, for all $m,r,t\in\mathbb{N}$ with $r\le m$ and $t\le m$,
$\bigl(3r+14T_c(r)+50\bigr)(t+1)^2\le A(m+1)^d$, where $T_c$ is the canonizer time
function of $c$.
\end{lemma}

\begin{definition}[Quadruple encoding]
\lean{Complexity.tmsatQuad}
\label{Complexity.tmsatQuad}
\leanok
\uses{Turing.pairEncode}
For bit strings $\alpha, x$ and $n, t \in \mathbb{N}$, the quadruple encoding $\langle\alpha, x, 1^n, 1^t\rangle$ is the bit string $\langle\alpha,\langle x,\langle 1^n,1^t\rangle\rangle\rangle$, where $\langle a,b\rangle$ is the self-delimiting pair encoding and $1^k$ is the word of $k$ ones.
\end{definition}

\begin{lemma}[Bounds on encoded component lengths]
\lean{Complexity.tmsat_quad_bounds}
\label{Complexity.tmsat_quad_bounds}
\leanok
\uses{Complexity.tmsatQuad, Turing.universal_pair_length}
\difficulty{2}
For all bit strings $\alpha,x$ and all natural numbers $n,t$, each of $|\alpha|$, $|x|$,
$n$, and $t$ is at most the length of the nested encoding $\langle\alpha,\langle
x,\langle 1^n,1^t\rangle\rangle\rangle$, where $1^j$ denotes the word of $j$ ones.
\end{lemma}

\begin{definition}[Verifier language for TMSAT]
\lean{Complexity.tmsatVerifier}
\label{Complexity.tmsatVerifier}
\leanok
\uses{Complexity.tmsatQuad, Turing.CodeTM.toFinTM, Turing.FinTM.ComputesInTime, Turing.MachineCode, Turing.pairEncode}
Let $c$ be a machine-representation scheme, which assigns to every bit string $\alpha$ a one-work-tape binary Turing machine $M_\alpha$ (decoding is total). The verifier language $V_c$ consists of all bit strings $yw$ such that $|w| = |y| + 1$, $y = \langle\alpha,\langle x,\langle 1^n,1^t\rangle\rangle\rangle$ for some bit strings $\alpha, x$ and some $n, t \in \mathbb{N}$ (self-delimiting pair encoding, $1^k$ the word of $k$ ones), and $M_\alpha$ on input $\langle x, w_{[:n]}\rangle$ halts within $t$ steps with output exactly $1$, where $w_{[:n]}$ consists of the first $n$ bits of $w$.
\end{definition}

\begin{lemma}[Length of a padded concatenation]
\lean{Complexity.tmsat_split_length}
\label{Complexity.tmsat_split_length}
\leanok
\difficulty{2}
Let $y$ and $w$ be finite strings of bits, and suppose $|w|=|y|+1$. Then $|yw|$ is odd,
and $(|yw|-1)/2=|y|$.
\end{lemma}

\begin{lemma}[TMSAT certificate characterization]
\lean{Complexity.tmsat_certificate_equiv}
\label{Complexity.tmsat_certificate_equiv}
\leanok
\uses{Complexity.TMSAT, Complexity.tmsatVerifier, Complexity.tmsat_quad_bounds, Turing.MachineCode}
\difficulty{5}
Let $c$ be a machine representation scheme and let $y$ be a bit string. Then
$y\\in\\mathrm{TMSAT}_c$ if and only if there is a bit string $w$ of length $|y|+1$ such
that $yw$ belongs to the verifier language $V_c$.
\end{lemma}

\begin{definition}[Pair-to-concatenation map]
\lean{Complexity.tmsatConcat}
\label{Complexity.tmsatConcat}
\leanok
\uses{Turing.pairDecode}
This is the function on bit strings that sends every valid self-delimiting pair encoding $\langle a,b\rangle$ (the bits of $a$ doubled, then $01$, then $b$) to the concatenation $ab$, and every bit string that does not parse as a pair encoding to the empty string.
\end{definition}

\begin{lemma}[Polynomial-time equality test]
\lean{Complexity.tmsat_pt_eq}
\label{Complexity.tmsat_pt_eq}
\leanok
\uses{Complexity.PolyTimeComputable, Complexity.polyTimeComputable_of_linear, Turing.FinTM.computesFunInTime_ifEq}
\difficulty{3}
For every fixed bit string $w$, the function that maps a bit string $x$ to the one-bit
word $[1]$ if $x=w$ and to the one-bit word $[0]$ otherwise is polynomial-time
computable.
\end{lemma}

\begin{lemma}[Nonempty fixed-width increment output]
\lean{Complexity.tmsat_inc_nonempty}
\label{Complexity.tmsat_inc_nonempty}
\leanok
\uses{Turing.incFixed}
\difficulty{3}
For every bit string $w$, if its fixed-width little-endian binary increment produces an
output, that output is not the empty word.
\end{lemma}

\begin{lemma}[Overflow of fixed-width increment]
\lean{Complexity.tmsat_inc_none}
\label{Complexity.tmsat_inc_none}
\leanok
\uses{Turing.incFixed}
\difficulty{3}
For every bit string $w$, the fixed-width little-endian binary increment overflows
exactly when every bit of $w$ is $1$, including the empty word as the all-ones word of
length zero.
\end{lemma}

\begin{lemma}[All-ones test in polynomial time]
\lean{Complexity.tmsat_pt_unary}
\label{Complexity.tmsat_pt_unary}
\leanok
\uses{Complexity.PolyTimeComputable, Complexity.PolyTimeComputable.comp, Complexity.polyTimeComputable_of_linear, Complexity.tmsat_inc_none, Complexity.tmsat_inc_nonempty, Complexity.tmsat_pt_eq, Turing.FinTM.computesFunInTime_incFixed, Turing.incFixed}
\difficulty{4}
The function sending a bit string $x$ to the one-bit word whose bit is $1$ exactly when
every bit of $x$ is $1$ is polynomial-time computable. The condition includes the empty
word, for which the output bit is $1$.
\end{lemma}

\begin{definition}[Prefix-extraction machine]
\lean{Complexity.tmsatTakeTM}
\label{Complexity.tmsatTakeTM}
\leanok
\uses{Turing.FinTM}
The prefix extractor $E$ is the finite binary-alphabet machine with one work tape, states $\{0,1,2,3\}$ and initial state $0$ which, on an input $\langle a,b\rangle$ (the bits of $a$ doubled, then $01$, then $b$), reads the doubled region in aligned blocks of two bits: each data block $00$ or $11$ writes one more $1$ on a unary counter on its work tape and moves the work head right, and the separator $01$ moves the work head left onto the last counter cell and enters state $3$. In state $3$ it copies the next input bit to the output and moves its input head right and its work head left, halting as soon as the work head reads a blank or the input is exhausted; a malformed block $10$, or the end of the input in states $0$ to $2$, halts it immediately without output.
\end{definition}

\begin{definition}[Configurations of the prefix extractor]
\lean{Complexity.tmsatTakeCfg}
\label{Complexity.tmsatTakeCfg}
\leanok
\uses{Complexity.polyTape, Turing.Cfg}
For a bit string $x$, an optional state $q \in \{0,1,2,3\}$ (absence meaning halted), an index $i \le |x|$, $m \in \mathbb{N}$, a work-head position $h \in \mathbb{Z}$ and an output word $o$, $\Delta_x(q,i,m,h,o)$ is the configuration of a one-work-tape binary machine with states $\{0,1,2,3\}$ on input $x$ in state $q$, with the input head on input position $i$ (counting from $0$; the blank cell right of the input when $i = |x|$), the work tape holding $1^m$ in cells $0,\dots,m-1$ and blanks elsewhere, the work head at $h$, and output $o$.
\end{definition}

\begin{lemma}[Input symbol at a configuration position]
\lean{Complexity.tmsat_take_read}
\label{Complexity.tmsat_take_read}
\leanok
\uses{Complexity.tmsatTakeCfg, Turing.Cfg.inputSymbol, Turing.FinTM.inputSymbol_at}
\difficulty{1}
Let $x$ be a bit string, let $q$ be an optional state from a four-state set, and let
$i,m\in\mathbb N$, $h\in\mathbb Z$, and $o$ be a bit string. Suppose $i\le |x|$. In the
configuration on input $x$ with state $q$, input head at position $i$, unary work tape
content $1^m$, work head at $h$, and output $o$, the symbol under the input head is the
bit at position $i$ when $i<|x|$, and is blank when $i=|x|$.
\end{lemma}

\begin{lemma}[Counting a doubled data block]
\lean{Complexity.tmsat_take_double}
\label{Complexity.tmsat_take_double}
\leanok
\uses{Complexity.polyTape_write, Complexity.tmsatTakeCfg, Complexity.tmsatTakeTM, Complexity.tmsat_take_read, Turing.Action, Turing.Action.apply, Turing.MultiTapeTM.runFrom, Turing.MultiTapeTM.runFrom_succ_eq_step, Turing.MultiTapeTM.runFrom_zero, Turing.MultiTapeTM.step, Turing.moveInputPos_pos_of_ne_right}
\difficulty{4}
Let $x$ be a bit string, let $i,m\in\mathbb N$, and let $b\in\{0,1\}$. Suppose $i+2\le
|x|$ and the bits of $x$ at positions $i$ and $i+1$ are both $b$. After two steps of the
prefix extractor from state $0$, with its input head at position $i$, its unary work
tape containing $1^m$, its work head at position $m$, and empty output, the machine is
in state $0$ with its input head at $i+2$, its work tape containing $1^{m+1}$, its work
head at $m+1$, and empty output.
\end{lemma}

\begin{lemma}[Reading the pair separator]
\lean{Complexity.tmsat_take_separator}
\label{Complexity.tmsat_take_separator}
\leanok
\uses{Complexity.tmsatTakeCfg, Complexity.tmsatTakeTM, Complexity.tmsat_take_read, Turing.Action, Turing.Action.apply, Turing.MultiTapeTM.runFrom, Turing.MultiTapeTM.runFrom_succ_eq_step, Turing.MultiTapeTM.runFrom_zero, Turing.MultiTapeTM.step, Turing.moveInputPos_pos_of_ne_right}
\difficulty{3}
Let $x$ be a bit string and $i,m\in\mathbb N$. Suppose $i+2\le |x|$, the bit at position
$i$ is $0$, and the bit at position $i+1$ is $1$. Starting from state $0$ with the input
head at $i$, unary work tape content $1^m$, work head at $m$, and empty output, after
two steps the extractor is in state $3$, with the input head at $i+2$, work tape content
$1^m$, work head at $m-1$, and empty output.
\end{lemma}

\begin{lemma}[Parsing the encoded first component]
\lean{Complexity.tmsat_take_parse}
\label{Complexity.tmsat_take_parse}
\leanok
\uses{Complexity.tmsatTakeCfg, Complexity.tmsatTakeTM, Complexity.tmsat_take_double, Complexity.tmsat_take_separator, Turing.MultiTapeTM.runFrom, Turing.MultiTapeTM.runFrom_add, Turing.pairEncode, Turing.universal_pair_length}
\difficulty{4}
Let $a,b,x$, and $p$ be bit strings, let $m\in\mathbb N$, and suppose $x=p\langle
a,b\rangle$, where the encoding doubles every bit of $a$, appends $01$, and then appends
$b$. Starting from state $0$ with the input head at position $|p|$, unary work tape
content $1^m$, work head at $m$, and empty output, after $2|a|+2$ steps the extractor is
in state $3$, with its input head at $|p|+2|a|+2$, work tape content $1^{m+|a|}$, work
head at $m+|a|-1$, and empty output.
\end{lemma}

\begin{lemma}[Copying a payload prefix]
\lean{Complexity.tmsat_take_payload}
\label{Complexity.tmsat_take_payload}
\leanok
\uses{Complexity.polyTape, Complexity.tmsatTakeCfg, Complexity.tmsatTakeTM, Complexity.tmsat_take_read, Turing.Action.apply, Turing.Cfg.workTapeSymbols, Turing.MultiTapeTM.runFrom, Turing.MultiTapeTM.runFrom_succ_eq_step', Turing.moveInputPos_pos_of_ne_right}
\difficulty{4}
Let $x,p,b$ be bit strings and $m,j\in\mathbb N$. Suppose $x=p b$ and $j\le\min(m,|b|)$.
Starting in state $3$ with the input head at $|p|$, unary work tape content $1^m$, work
head at $m-1$, and empty output, after $j$ steps the extractor is in state $3$, with
input head at $|p|+j$, the same work tape content, work head at $m-j-1$, and output
equal to the first $j$ bits of $b$.
\end{lemma}

\begin{lemma}[Prefix extraction from a pair encoding]
\lean{Complexity.tmsat_take_computes}
\label{Complexity.tmsat_take_computes}
\leanok
\uses{Complexity.polyTape, Complexity.tmsatTakeCfg, Complexity.tmsatTakeTM, Complexity.tmsat_take_parse, Complexity.tmsat_take_payload, Complexity.tmsat_take_read, Turing.Action.apply, Turing.Cfg.workTapeSymbols, Turing.FinTM.ComputesInTime, Turing.FinTM.ComputesInTime.mono, Turing.FinTM.computesInTime_iff, Turing.MultiTapeTM.initCfg, Turing.MultiTapeTM.runFrom, Turing.MultiTapeTM.runFrom_add, Turing.MultiTapeTM.runFrom_succ_eq_step', Turing.pairEncode, Turing.universal_pair_length}
\difficulty{5}
For any bit strings $a$ and $b$, the prefix extractor, started on the encoding that
doubles each bit of $a$, appends the separator $01$, and then appends $b$, halts within
$|\langle a,b\rangle|+1$ steps and outputs the first $|a|$ bits of $b$, or all of $b$ if
$|b|<|a|$.
\end{lemma}

\begin{lemma}[Truncation to a computed length]
\lean{Complexity.tmsat_pt_take}
\label{Complexity.tmsat_pt_take}
\leanok
\uses{Complexity.PolyTimeComputable, Complexity.PolyTimeComputable.pairEncode, Complexity.tmsatTakeTM, Complexity.tmsat_take_computes, Turing.FinTM.ComputesInTime, Turing.FinTM.ComputesInTime.mono, Turing.FinTM.computesInTime_iff, Turing.FinTM.exists_comp_on_image, Turing.MultiTapeTM.output_length_le, Turing.pairEncode}
\difficulty{4}
If $f,g : \{0,1\}^* \to \{0,1\}^*$ are polynomial-time computable, then the function
sending each bit string $x$ to the prefix of $g(x)$ of length $|f(x)|$ is
polynomial-time computable; if $g(x)$ is shorter, this prefix is all of $g(x)$.
\end{lemma}

\begin{lemma}[Reconstruction of a parsed pair]
\lean{Complexity.tmsat_pair_valid}
\label{Complexity.tmsat_pair_valid}
\leanok
\uses{Complexity.pairFstD, Complexity.pairSndD, Turing.eq_pairEncode_of_pairDecode, Turing.pairDecode, Turing.pairEncode}
\difficulty{3}
Let $z$ be a bit string for which the aligned pair parser succeeds, and let $\pi_1(z)$
and $\pi_2(z)$ be its decoded components. Then $z=\langle\pi_1(z),\pi_2(z)\rangle$,
where the pairing doubles each bit of the first component, appends the separator $01$,
and then appends the second component.
\end{lemma}

\begin{definition}[Odd split of a string]
\lean{Complexity.tmsatSplit}
\label{Complexity.tmsatSplit}
\leanok
\uses{Turing.pairEncode, Turing.solveSplit}
For a bit string $z$, the odd split $s(z)$ is the self-delimiting pair encoding $\langle z_{[:i]}, z_{[i:]}\rangle$ of the first $i$ bits of $z$ and the remaining bits, where $i$ is the least index in $\{0,\dots,|z|\}$ with $i + (i+1) = |z|$; if no such index exists (that is, $|z|$ is even), $s(z)$ is the empty string.
\end{definition}

\begin{lemma}[Equation satisfied by a successful split]
\lean{Complexity.tmsat_split_some}
\label{Complexity.tmsat_split_some}
\leanok
\uses{Turing.solveSplit}
\difficulty{2}
For $N\in\mathbb N$, if the bounded search for the least $i\in\{0,\ldots,N\}$ satisfying
$i+(i+1)=N$ returns $i$, then $i+(i+1)=N$.
\end{lemma}

\begin{lemma}[Split search finds the solution]
\lean{Complexity.tmsat_split_exists}
\label{Complexity.tmsat_split_exists}
\leanok
\uses{Complexity.tmsat_split_some, Turing.solveSplit}
\difficulty{3}
For natural numbers $N$ and $i$ satisfying $i+(i+1)=N$, the bounded search for the least
$j\in\{0,\ldots,N\}$ satisfying $j+(j+1)=N$ returns $i$.
\end{lemma}

\begin{lemma}[Split of a padded string]
\lean{Complexity.tmsat_split_append}
\label{Complexity.tmsat_split_append}
\leanok
\uses{Complexity.tmsatSplit, Complexity.tmsat_split_exists, Turing.pairEncode}
\difficulty{3}
Let $y,w$ be bit strings with $|w|=|y|+1$. The odd split of their concatenation $yw$ is
the self-delimiting pair encoding $\langle y,w\rangle$.
\end{lemma}

\begin{lemma}[Components of a valid odd split]
\lean{Complexity.tmsat_split_components}
\label{Complexity.tmsat_split_components}
\leanok
\uses{Complexity.pairFstD, Complexity.pairSndD, Complexity.tmsatSplit, Complexity.tmsat_split_some, Turing.pairDecode, Turing.pairDecode_pairEncode, Turing.solveSplit}
\difficulty{4}
Let $z$ be a bit string whose odd split parses as a pair encoding. The total first and
second projections of that split concatenate to $z$, and the second projected component
has length one greater than the first.
\end{lemma}

\begin{definition}[Instance part of a verifier input]
\lean{Complexity.tmsatY}
\label{Complexity.tmsatY}
\leanok
\uses{Complexity.pairFstD, Complexity.tmsatSplit}
For a bit string $z$, the instance part is $Y(z) = \pi_1(s(z))$, where $s(z)$ is the odd split of $z$ (the self-delimiting pair encoding $\langle z_{[:i]}, z_{[i:]}\rangle$ of the first $i$ bits of $z$ and the rest when $|z| = 2i+1$, and the empty word when $|z|$ is even), and $\pi_1, \pi_2$ return the first and second components of a pair encoding (the empty word on strings that do not parse as pairs).
\end{definition}

\begin{definition}[Certificate part of a verifier input]
\lean{Complexity.tmsatW}
\label{Complexity.tmsatW}
\leanok
\uses{Complexity.pairSndD, Complexity.tmsatSplit}
For a bit string $z$, the padded certificate part is $W(z) = \pi_2(s(z))$, where $s(z)$ is the odd split of $z$ (the self-delimiting pair encoding $\langle z_{[:i]}, z_{[i:]}\rangle$ of the first $i$ bits of $z$ and the rest when $|z| = 2i+1$, and the empty word when $|z|$ is even), and $\pi_1, \pi_2$ return the first and second components of a pair encoding (the empty word on strings that do not parse as pairs).
\end{definition}

\begin{definition}[Code field of a verifier input]
\lean{Complexity.tmsatCode}
\label{Complexity.tmsatCode}
\leanok
\uses{Complexity.pairFstD, Complexity.tmsatY}
For a bit string $z$, the code field is $\pi_1(Y(z))$, where $Y(z) = \pi_1(s(z))$ and $s(z)$ is the odd split of $z$ (the self-delimiting pair encoding $\langle z_{[:i]}, z_{[i:]}\rangle$ of the first $i$ bits of $z$ and the rest when $|z| = 2i+1$, and the empty word when $|z|$ is even), and $\pi_1, \pi_2$ return the first and second components of a pair encoding (the empty word on strings that do not parse as pairs).
\end{definition}

\begin{definition}[Input field of a verifier input]
\lean{Complexity.tmsatInput}
\label{Complexity.tmsatInput}
\leanok
\uses{Complexity.pairFstD, Complexity.pairSndD, Complexity.tmsatY}
For a bit string $z$, the input field is $\pi_1(\pi_2(Y(z)))$, where $Y(z) = \pi_1(s(z))$ and $s(z)$ is the odd split of $z$ (the self-delimiting pair encoding $\langle z_{[:i]}, z_{[i:]}\rangle$ of the first $i$ bits of $z$ and the rest when $|z| = 2i+1$, and the empty word when $|z|$ is even), and $\pi_1, \pi_2$ return the first and second components of a pair encoding (the empty word on strings that do not parse as pairs).
\end{definition}

\begin{definition}[Certificate-width field of a verifier input]
\lean{Complexity.tmsatWidth}
\label{Complexity.tmsatWidth}
\leanok
\uses{Complexity.pairFstD, Complexity.pairSndD, Complexity.tmsatY}
For a bit string $z$, the width field is $\pi_1(\pi_2(\pi_2(Y(z))))$, where $Y(z) = \pi_1(s(z))$ and $s(z)$ is the odd split of $z$ (the self-delimiting pair encoding $\langle z_{[:i]}, z_{[i:]}\rangle$ of the first $i$ bits of $z$ and the rest when $|z| = 2i+1$, and the empty word when $|z|$ is even), and $\pi_1, \pi_2$ return the first and second components of a pair encoding (the empty word on strings that do not parse as pairs); on well-formed inputs it is the unary certificate length $1^n$, but its shape is not checked.
\end{definition}

\begin{definition}[Deadline field of a verifier input]
\lean{Complexity.tmsatClock}
\label{Complexity.tmsatClock}
\leanok
\uses{Complexity.pairSndD, Complexity.tmsatY}
For a bit string $z$, the clock field is $\pi_2(\pi_2(\pi_2(Y(z))))$, where $Y(z) = \pi_1(s(z))$ and $s(z)$ is the odd split of $z$ (the self-delimiting pair encoding $\langle z_{[:i]}, z_{[i:]}\rangle$ of the first $i$ bits of $z$ and the rest when $|z| = 2i+1$, and the empty word when $|z|$ is even), and $\pi_1, \pi_2$ return the first and second components of a pair encoding (the empty word on strings that do not parse as pairs); on well-formed inputs it is the unary deadline $1^t$, but its shape is not checked.
\end{definition}

\begin{definition}[Validity guard for verifier inputs]
\lean{Complexity.tmsatGood}
\label{Complexity.tmsatGood}
\leanok
\uses{Complexity.pairSndD, Complexity.tmsatClock, Complexity.tmsatSplit, Complexity.tmsatWidth, Complexity.tmsatY, Turing.pairDecode}
For a bit string $z$, write $s(z)$ for its odd split (the self-delimiting pair encoding $\langle z_{[:i]}, z_{[i:]}\rangle$ when $|z| = 2i+1$, and the empty word when $|z|$ is even), $\pi_1, \pi_2$ for the two components of a pair encoding (the empty word when parsing fails), $Y = \pi_1(s(z))$ and $W = \pi_2(s(z))$; the code, input, width and clock fields of $z$ are $\pi_1(Y)$, $\pi_1(\pi_2(Y))$, $\pi_1(\pi_2(\pi_2(Y)))$ and $\pi_2(\pi_2(\pi_2(Y)))$. The guard $G(z)$ is the Boolean that is true exactly when $s(z)$, $Y$, $\pi_2(Y)$ and $\pi_2(\pi_2(Y))$ all parse as pair encodings and both the width field and the clock field consist only of ones (the empty word included).
\end{definition}

\begin{lemma}[Structure of valid verifier inputs]
\lean{Complexity.tmsat_good_spec}
\label{Complexity.tmsat_good_spec}
\leanok
\uses{Complexity.pairSndD, Complexity.tmsatClock, Complexity.tmsatCode, Complexity.tmsatGood, Complexity.tmsatInput, Complexity.tmsatQuad, Complexity.tmsatW, Complexity.tmsatWidth, Complexity.tmsatY, Complexity.tmsat_pair_valid, Complexity.tmsat_split_components, Turing.pairEncode}
\difficulty{4}
Let $z$ be a bit string satisfying the validity guard. Write $Y$ for its instance part
and $W$ for its certificate part; write $\alpha$, $x$, $1^n$, and $1^t$ for the code,
input, width, and clock fields extracted from $Y$, respectively, with $n$ and $t$ the
lengths of the width and clock fields. Then $z=YW$, the length of $W$ is one more than
the length of $Y$, and $Y=\langle\alpha,\langle x,\langle1^n,1^t\rangle\rangle\rangle$.
\end{lemma}

\begin{lemma}[Fields of a padded quadruple]
\lean{Complexity.tmsat_good_quad}
\label{Complexity.tmsat_good_quad}
\leanok
\uses{Complexity.pairFstD, Complexity.pairSndD, Complexity.tmsatClock, Complexity.tmsatCode, Complexity.tmsatGood, Complexity.tmsatInput, Complexity.tmsatQuad, Complexity.tmsatW, Complexity.tmsatWidth, Complexity.tmsatY, Complexity.tmsat_split_append, Turing.pairDecode_pairEncode}
\difficulty{3}
Let $\alpha$, $x$, and $w$ be bit strings, and let $n,t\in\mathbb{N}$. Suppose the
length of $w$ is one more than the length of the encoding
$\langle\alpha,x,1^n,1^t\rangle$, where $1^k$ denotes the word of $k$ ones. Append $w$
to this encoding to form $z$. Then $z$ satisfies the validity guard; its code field is
$\alpha$, its input field is $x$, its width and clock fields have lengths $n$ and $t$,
respectively, and its certificate part is $w$.
\end{lemma}

\begin{lemma}[Polynomial-time verifier fields]
\lean{Complexity.tmsat_fields_poly}
\label{Complexity.tmsat_fields_poly}
\leanok
\uses{Complexity.PolyTimeComputable, Complexity.PolyTimeComputable.comp, Complexity.pairFstD, Complexity.pairSndD, Complexity.polyTimeComputable_and, Complexity.polyTimeComputable_of_linear, Complexity.polyTimeComputable_pairFstD, Complexity.polyTimeComputable_pairSndD, Complexity.tmsatClock, Complexity.tmsatCode, Complexity.tmsatGood, Complexity.tmsatInput, Complexity.tmsatSplit, Complexity.tmsatW, Complexity.tmsatWidth, Complexity.tmsatY, Complexity.tmsat_pt_unary, Turing.FinTM.computesFunInTime_pairValid, Turing.FinTM.computesFunInTime_splitSolve}
\difficulty{4}
For every bit string $z$, each of the following functions is polynomial-time computable:
its code field, input field, unary width field, unary clock field, and certificate
component, as obtained by the total parsing and projection operations defining these
fields; and the function returning the one-bit word encoding whether the specified
parsing and unary-field validity checks all hold.
\end{lemma}

\begin{definition}[Guarded timed simulation request]
\lean{Complexity.tmsatRequest}
\label{Complexity.tmsatRequest}
\leanok
\uses{Complexity.tmsatClock, Complexity.tmsatCode, Complexity.tmsatGood, Complexity.tmsatInput, Complexity.tmsatW, Complexity.tmsatWidth, Turing.pairEncode}
For a bit string $z$, write $s(z)$ for its odd split (the self-delimiting pair encoding $\langle z_{[:i]}, z_{[i:]}\rangle$ when $|z| = 2i+1$, and the empty word when $|z|$ is even), $\pi_1, \pi_2$ for the two components of a pair encoding (the empty word when parsing fails), $Y = \pi_1(s(z))$ and $W = \pi_2(s(z))$; the code, input, width and clock fields of $z$ are $\pi_1(Y)$, $\pi_1(\pi_2(Y))$, $\pi_1(\pi_2(\pi_2(Y)))$ and $\pi_2(\pi_2(\pi_2(Y)))$. The guard $G(z)$ holds when $s(z)$, $Y$, $\pi_2(Y)$ and $\pi_2(\pi_2(Y))$ all parse as pair encodings and the width and clock fields consist only of ones. Then the timed request of $z$ is $\langle\langle \mathrm{bin}(t), \alpha\rangle, \langle x, W_{[:n]}\rangle\rangle$ when $G(z)$ holds, where $\alpha$ and $x$ are the code and input fields, $n$ and $t$ the lengths of the width and clock fields, $W_{[:n]}$ the first $n$ bits of $W$, and $\mathrm{bin}(t)$ the little-endian binary representation of $t$; otherwise it is the fixed request $\langle\langle\varepsilon,\varepsilon\rangle,\varepsilon\rangle$.
\end{definition}

\begin{definition}[Guarded simulation result]
\lean{Complexity.tmsatResult}
\label{Complexity.tmsatResult}
\leanok
\uses{Complexity.tmsatAnswer, Complexity.tmsatClock, Complexity.tmsatCode, Complexity.tmsatGood, Complexity.tmsatInput, Complexity.tmsatW, Complexity.tmsatWidth, Turing.MachineCode, Turing.pairEncode}
Let $c$ be a machine-representation scheme, which assigns to every bit string $\alpha$ a one-work-tape binary Turing machine $M_\alpha$ (decoding is total). For a bit string $z$, write $s(z)$ for its odd split (the self-delimiting pair encoding $\langle z_{[:i]}, z_{[i:]}\rangle$ when $|z| = 2i+1$, and the empty word when $|z|$ is even), $\pi_1, \pi_2$ for the two components of a pair encoding (the empty word when parsing fails), $Y = \pi_1(s(z))$ and $W = \pi_2(s(z))$; the code, input, width and clock fields of $z$ are $\pi_1(Y)$, $\pi_1(\pi_2(Y))$, $\pi_1(\pi_2(\pi_2(Y)))$ and $\pi_2(\pi_2(\pi_2(Y)))$. The guard $G(z)$ holds when $s(z)$, $Y$, $\pi_2(Y)$ and $\pi_2(\pi_2(Y))$ all parse as pair encodings and the width and clock fields consist only of ones. Then the result for $z$ is $A_c(\alpha, \langle x, W_{[:n]}\rangle, t)$ when $G(z)$ holds, with $\alpha$, $x$ the code and input fields, $n$, $t$ the lengths of the width and clock fields and $W_{[:n]}$ the first $n$ bits of $W$, and $A_c(\varepsilon,\varepsilon,0)$ otherwise; here $A_c(\beta, v, t')$ is $1w'$ if $M_\beta$ run on $v$ for $t'$ steps has halted with output $w'$, and the single bit $0$ if it has not halted.
\end{definition}

\begin{lemma}[Polynomial-time request assembly]
\lean{Complexity.tmsat_request_poly}
\label{Complexity.tmsat_request_poly}
\leanok
\uses{Complexity.PolyTimeComputable, Complexity.PolyTimeComputable.comp, Complexity.PolyTimeComputable.pairEncode, Complexity.PolyTimeComputable.pairMapSnd, Complexity.pairMapSnd, Complexity.polyTimeComputable_const, Complexity.polyTimeComputable_id, Complexity.polyTimeComputable_ite, Complexity.polyTimeComputable_of_linear, Complexity.tmsatClock, Complexity.tmsatRequest, Complexity.tmsat_fields_poly, Complexity.tmsat_pt_take, Turing.FinTM.computesFunInTime_lengthBits, Turing.FinTM.computesFunInTime_pairSnd, Turing.pairDecode_pairEncode, Turing.pairEncode}
\difficulty{4}
The map sending a bit string $z$ to its timed request is polynomial-time computable. To
define it, split $z$ at the least $i$ satisfying $i+(i+1)=|z|$ and pair-encode the two
resulting parts; if there is no such $i$, use the empty word. Parse this encoding and
its nested fields using the total pair projections, which return the empty word when
parsing fails. If all required pair parses succeed and the width and clock fields
consist only of ones, output the pair encoding whose first component is the pair of the
little-endian binary representation of the clock-field length and the code field, and
whose second component is the pair of the input field and the prefix of the padded
certificate field of length equal to the width-field length; otherwise output the fixed
encoding $\langle\langle\varepsilon,\varepsilon\rangle,\varepsilon\rangle$.
\end{lemma}

\begin{lemma}[Bounds on valid code and clock fields]
\lean{Complexity.tmsat_good_bounds}
\label{Complexity.tmsat_good_bounds}
\leanok
\uses{Complexity.tmsatClock, Complexity.tmsatCode, Complexity.tmsatGood, Complexity.tmsatInput, Complexity.tmsatWidth, Complexity.tmsat_good_spec, Complexity.tmsat_quad_bounds}
\difficulty{3}
For a bit string $z$ satisfying the validity guard, the code field and clock field each
have length at most the length of $z$.
\end{lemma}

\begin{lemma}[Acceptance characterizes the verifier language]
\lean{Complexity.tmsat_result_accept}
\label{Complexity.tmsat_result_accept}
\leanok
\uses{Complexity.tmsatAnswer_accept, Complexity.tmsatClock, Complexity.tmsatCode, Complexity.tmsatGood, Complexity.tmsatInput, Complexity.tmsatResult, Complexity.tmsatVerifier, Complexity.tmsatW, Complexity.tmsatWidth, Complexity.tmsatY, Complexity.tmsat_good_quad, Complexity.tmsat_good_spec, Turing.CodeTM.toFinTM, Turing.FinTM.ComputesInTime, Turing.FinTM.not_computesInTime_zero, Turing.MachineCode}
\difficulty{4}
Let $c$ be a machine representation scheme and $z$ a bit string. The result produced
from $z$ by the guarded computation equals the two-bit word $11$ if and only if $z$
belongs to the verifier language: the set of words $y'w'$ such that $|w'|=|y'|+1$ and
there exist bit strings $\beta,v$ and natural numbers $n',t'$ with
$y'=\langle\beta,\langle v,\langle1^{n'},1^{t'}\rangle\rangle\rangle$, for which the
machine decoded from $\beta$ halts within $t'$ steps on input $\langle v,w'_{[:n']
}\rangle$ with output exactly $1$.
\end{lemma}

\begin{theorem}[NP membership of bounded machine acceptance]
\lean{Complexity.TMSAT_mem_NP}
\label{Complexity.TMSAT_mem_NP}
\leanok
\uses{Complexity.NP, Complexity.P, Complexity.PolyBound, Complexity.PolyTimeComputable, Complexity.PolyTimeComputable.comp, Complexity.TMSAT, Complexity.mem_P_iff, Complexity.timed_universal_quantitative, Complexity.tmsatClock, Complexity.tmsatCode, Complexity.tmsatGood, Complexity.tmsatInput, Complexity.tmsatRequest, Complexity.tmsatResult, Complexity.tmsatVerifier, Complexity.tmsatW, Complexity.tmsatWidth, Complexity.tmsat_certificate_equiv, Complexity.tmsat_good_bounds, Complexity.tmsat_pt_eq, Complexity.tmsat_request_poly, Complexity.tmsat_result_accept, Complexity.tmsat_simulation_budget, Complexity.tmsat_simulator_total, Turing.EffectiveMachineCode, Turing.FinTM.ComputesInTime, Turing.FinTM.ComputesInTime.mono, Turing.FinTM.exists_comp_on_image, Turing.MultiTapeTM.indicator, Turing.pairEncode}
\difficulty{5}
Let $c$ be an effective machine-code scheme whose canonizer has a polynomially bounded
running-time function. The language of strings that can be written as
$\langle\alpha,\langle x,\langle 1^n,1^t\rangle\rangle$, where $1^m$ denotes a string of
$m$ ones, and for which there is a string $u$ of length $n$ such that the machine
decoded from $\alpha$ halts after exactly $t$ steps on input $\langle x,u\rangle$ and
outputs the single bit $1$, belongs to NP.
\end{theorem}

\begin{lemma}[Arithmetic deadline bound]
\lean{Complexity.tmsat_deadline_bound}
\label{Complexity.tmsat_deadline_bound}
\leanok
\difficulty{4}
For all $K,B,e,C,c,n\\in\\mathbb{N}$, $K\\bigl(B(2n+2+C(n+1)^c+1)^e+1\\bigr)^2\\le
((K+1)(B+1)^2(C+3)^{2e})(n+1)^{2e\\max(1,c)}$.
\end{lemma}

\begin{lemma}[Constant functions are polynomial-time]
\lean{Complexity.tmsat_constant_poly}
\label{Complexity.tmsat_constant_poly}
\leanok
\uses{Complexity.PolyTimeComputable, Turing.FinTM.computesFunInTime_const}
\difficulty{2}
For every bit string $w$, the function sending every bit string to $w$ is
polynomial-time computable.
\end{lemma}

\begin{lemma}[Polynomial-time binary certificate length]
\lean{Complexity.tmsat_exact_certificate_bits}
\label{Complexity.tmsat_exact_certificate_bits}
\leanok
\uses{Complexity.PolyTimeComputable, Complexity.timeConstructible_poly, Complexity.tmsat_constant_poly, Turing.FinTM.ComputesInTime.mono}
\difficulty{4}
For all $C,c\\in\\mathbb{N}$, the function sending a bit string $x$ to the bit-list
representation of the natural number $C(|x|+1)^c$ is polynomial-time computable.
\end{lemma}

\begin{definition}[Output function of the verifier wrapper]
\lean{Complexity.tmsatWrapperOutput}
\label{Complexity.tmsatWrapperOutput}
\leanok
\uses{Turing.MultiTapeTM.indicator, Turing.pairDecode}
For a language $V$ of bit strings, this is the total function on bit strings that sends a valid self-delimiting pair encoding $\langle x,u\rangle$ to the single bit $1$ if the concatenation $xu$ lies in $V$ and to $0$ otherwise, and sends every string that does not parse as a pair encoding to the single bit $0$.
\end{definition}

\begin{lemma}[Injectivity of the quadruple encoding]
\lean{Complexity.tmsat_quad_injective}
\label{Complexity.tmsat_quad_injective}
\leanok
\uses{Complexity.tmsatQuad, Turing.pairEncode, Turing.pairEncode_injective}
\difficulty{3}
For bit strings $\alpha,x,\beta,y$ and natural numbers $n,t,m,s$, if
$\langle\alpha,\langle x,\langle 1^n,1^t\rangle\rangle\rangle=\langle\beta,\langle
y,\langle 1^m,1^s\rangle\rangle\rangle$, then $\alpha=\beta$, $x=y$, $n=m$, and $t=s$.
\end{lemma}

\begin{lemma}[Correctness of the TMSAT reduction]
\lean{Complexity.tmsat_reduction_correct}
\label{Complexity.tmsat_reduction_correct}
\leanok
\uses{Complexity.TMSAT, Complexity.tmsatQuad, Complexity.tmsat_quad_injective, Turing.CodeTM.toFinTM, Turing.FinTM.ComputesInTime, Turing.FinTM.ComputesInTime.output_unique, Turing.MachineCode, Turing.MultiTapeTM.indicator, Turing.pairEncode}
\difficulty{4}
Let $c$ be a machine representation scheme, let $L,V$ be languages of bit strings, let
$C,e$ be natural numbers, let $\alpha$ be a bit string, and let $T :
\mathbb{N}\to\mathbb{N}$. Suppose that for every bit string $x$, membership $x\in L$ is
equivalent to the existence of a bit string $u$ with $|u|=C(|x|+1)^e$ and $x u\in V$.
Suppose also that for every such $x,u$, the machine decoded from $\alpha$ halts within
$T(|x|)$ steps on input $\langle x,u\rangle$ and outputs the single bit indicating
whether $x u\in V$. Then for every bit string $x$, $x\in L$ if and only if
$\langle\alpha,\langle x,\langle1^{C(|x|+1)^e},1^{T(|x|)}\rangle\rangle\rangle$ belongs
to $\mathrm{TMSAT}_c$.
\end{lemma}

\begin{lemma}[Unary certificate length computation]
\lean{Complexity.tmsat_certificate_unary}
\label{Complexity.tmsat_certificate_unary}
\leanok
\uses{Complexity.PolyTimeComputable, Complexity.polyTimeComputable_const, Turing.FinTM.computesFunInTime_polyUnary}
\difficulty{3}
For all $C,c\\in\\mathbb{N}$, the function on bit strings sending $x$ to the word
consisting of $C(|x|+1)^c$ copies of $1$ is polynomial-time computable.
\end{lemma}

\begin{theorem}[TMSAT is NP-hard]
\lean{Complexity.TMSAT_NPHard}
\label{Complexity.TMSAT_NPHard}
\leanok
\uses{Complexity.NPHard, Complexity.PolyTimeComputable, Complexity.PolyTimeComputable.comp, Complexity.PolyTimeComputable.pairEncode, Complexity.TMSAT, Complexity.TimeConstructible, Complexity.mem_P_iff, Complexity.polyTimeComputable_const, Complexity.polyTimeComputable_id, Complexity.polyTimeComputable_ite, Complexity.polyTimeComputable_of_linear, Complexity.polyUnaryTM, Complexity.poly_unary_computes, Complexity.timeConstructible_poly, Complexity.tmsatConcat, Complexity.tmsatQuad, Complexity.tmsatWrapperOutput, Complexity.tmsat_certificate_unary, Complexity.tmsat_deadline_bound, Complexity.tmsat_exact_certificate_bits, Complexity.tmsat_reduction_correct, Turing.CodeTM.toFinTM, Turing.FinTM, Turing.FinTM.ComputesFunInTime, Turing.FinTM.ComputesInTime, Turing.FinTM.ComputesInTime.mono, Turing.FinTM.computesFunInTime_pairConcat, Turing.FinTM.computesFunInTime_pairEncodeFixed, Turing.FinTM.computesFunInTime_pairValid, Turing.FinTM.one_work_tape_binary, Turing.MachineCode, Turing.MachineCode.decode_encode, Turing.MultiTapeTM.indicator, Turing.exists_codeTM, Turing.pairDecode, Turing.pairDecode_pairEncode, Turing.pairEncode, Turing.universal_pair_length}
\difficulty{6}
For every machine-code scheme with a total decoder and the property that encoding any
machine followed by any number of $1$-bits decodes to that machine, the associated
language of encoded machine descriptions, inputs, unary witness-length bounds, and unary
time bounds for which some witness causes the decoded machine to halt with output $1$ is
NP-hard under polynomial-time many-one reductions.
\end{theorem}

\begin{theorem}[NP completeness of bounded machine acceptance]
\lean{Complexity.TMSAT_NPComplete}
\label{Complexity.TMSAT_NPComplete}
\leanok
\uses{Complexity.NPComplete, Complexity.PolyBound, Complexity.TMSAT, Complexity.TMSAT_NPHard, Complexity.TMSAT_mem_NP, Turing.EffectiveMachineCode}
\difficulty{1}
Let $c$ be an effective machine-code scheme whose canonizer has a polynomially bounded
running-time function. The language of strings that can be written as
$\langle\alpha,\langle x,\langle 1^n,1^t\rangle\rangle$, where $1^m$ denotes a string of
$m$ ones, and for which there is a string $u$ of length $n$ such that the machine
decoded from $\alpha$ halts after exactly $t$ steps on input $\langle x,u\rangle$ and
outputs the single bit $1$, is NP-complete.
\end{theorem}
